\subsection{How Do I Recover After a Failed Day?}

If the day failed, the first business goal is not to repair the whole self. It
is to stop the failed frame from reproducing itself tomorrow. Recovery works
best when you treat the failure as a state problem: identify the current state,
lower the barrier to the next usable target state, and spend only enough
volitional energy to re-enter motion inside the next decision window.

\begin{enumerate}[nosep]
  \item Name the failure in operational terms. What target state did you miss,
        and what current state won instead?
  \item Identify the main barrier that made the miss durable. Was it reward
        capture, ambiguity, environment, fatigue, or simple loss of timing?
  \item Choose a smaller target state for the next move. Do not try to recover
        the whole plan at once; choose the cheapest action that restores
        contact with the framework.
  \item Remove one source of friction before sleeping or stopping. Place the
        document, clothes, tools, or reminder where the next decision window
        can open with less resistance.
  \item Write the first visible action for tomorrow. The action should be so
        small that it can survive low volitional energy.
  \item Mark the next check point. Recovery is not complete when you feel
        better; it is complete when the next transition becomes easier to
        start.
\end{enumerate}

\begin{remark}
A failed day often looks like a moral collapse, but mechanically it is usually
just a strengthened inertia pattern. The best response is not a long apology to
the past day. It is a short intervention that makes the next target state more
available than the current one.
\end{remark}
